%% generated by mpce_results.py -- do not edit by hand
\begin{table}[!t]
\centering
\caption{Effect of the PSO Setting on the Comparison of the Previous Study's Methods (68 Cases, 6,030 Evaluations): Average Rank (1 = Best) with the Old ($w=0.7$, $c_1=c_2=2$) and the Constriction Setting, and Run-Level W/T/L of the Salp-Swarm Hybrids against PSO}
\label{tab:baseline}
\footnotesize\setlength{\tabcolsep}{4pt}
\begin{tabular}{lcc}
\toprule
& \multicolumn{2}{c}{PSO setting} \\
\cmidrule(lr){2-3}
Method & Old & Constriction \\
\midrule
LX-SSA-VNS & 2.65 & 3.41 \\
SSA-VNS & \textbf{1.99} & 2.77 \\
LX-SSA & 5.65 & 6.06 \\
SSA & 4.81 & 5.34 \\
PSO & 5.04 & \textbf{1.53} \\
DE & 6.99 & 7.03 \\
VNS & 3.59 & 4.15 \\
MS-SLSQP & 5.28 & 5.71 \\
\midrule
PSO position & 5 & 1 \\
LX-SSA-VNS vs.\ PSO (W/T/L) & 35/33/0 & 0/21/47 \\
SSA-VNS vs.\ PSO (W/T/L) & 41/27/0 & 0/22/46 \\
\bottomrule
\multicolumn{3}{p{0.95\columnwidth}}{Ranks: feasibility-aware rule (fewer than 15 of 30 feasible runs in a case = ranked last); bold: best average rank. W/T/L: cases in which the hybrid is significantly better / not different / worse than PSO (two-sided Wilcoxon signed-rank test, 30 seed-paired runs, Holm-adjusted over the seven comparisons of the hybrid in each case, $\alpha=0.05$).}
\end{tabular}
\end{table}
